% ======================================================================
%  Section 4.7.2 --- Convergence under occlusion.  COMPLETE.
%  All three occluded runs at iteration 600. Nothing is outstanding.
%
%  --------------------------------------------------------------------
%  1. THE RESULT CHANGED BETWEEN 400 AND 600, AND YOU NEED TO KNOW HOW.
%
%     At iteration 400 the two weighted variants looked like a clean tie:
%         V2  2.909 um / 0.01643 deg
%         V3  3.139 um / 0.01520 deg
%     7.7 per cent apart in translation and 7.4 per cent in rotation, in
%     opposite directions and by the same margin.
%
%     At iteration 600 they are not:
%         V2  2.474 um / 0.01646 deg
%         V3  3.156 um / 0.01519 deg
%     Variant 2 improved its translation by 15 per cent while Variant 3
%     did not move -- 3.139 to 3.156 um, marginally WORSE. So the gap
%     widened: Variant 2 is now better in translation by 27.6 per cent,
%     and Variant 3 remains better in rotation by 7.7 per cent.
%
%     This is a worse outcome for Variant 3 than the 400 reading suggested
%     and the section says so. It is still not a defeat, for a reason that
%     is quantitative and not a consolation: both weighted variants land
%     within three micrometres of the desired pose, against 1293 for the
%     baseline. The difference between 2.5 and 3.2 micrometres is of no
%     engineering consequence and cannot be established on one run each;
%     the difference between 3.2 and 1293 is the chapter's result.
%     What must go is any suggestion that Variant 3 positions better. It
%     does not.
%
%  --------------------------------------------------------------------
%  2. AND HERE IS THE FINDING THAT EXPLAINS IT.
%     V3's commanded velocity at 600 is 6.26e-6, against 1.83e-7 for V2 --
%     THIRTY-FOUR TIMES larger -- while its pose is unchanged from
%     iteration 400. The velocity is not unexhausted descent, as I
%     suggested last time it might be; it produces no net progress. That
%     is the other branch of the two I set out: Variant 3 carries a larger
%     terminal jitter, which is what one expects of weights recomputed
%     from the current image at every iteration, and it is the observable
%     trace of the time-varying weighting that 4.5 discusses without
%     measuring. I have written it in as an observation with that
%     connection, and NOT as a claim about the cause, which a single run
%     cannot support.
%     Consequent edit to 4.5 flagged at the foot.
%
%  --------------------------------------------------------------------
%  3. THE COMPUTATIONAL-COST WITHDRAWAL IS NOW BEYOND QUESTION.
%         iteration 400 : V1 8600 ms   V2 9227 ms   V3 11070 ms
%         iteration 600 : V1  591 ms   V2  496 ms   V3   406 ms
%     The 600 ordering is the exact REVERSE of the algorithmic one:
%     Variant 3 does the most work of the three and is the fastest of the
%     three, and the same code runs fourteen to twenty-seven times faster
%     than in the earlier snapshot. These timings measure machine load,
%     nothing else. The paragraph I wrote from the 400 figures, and its
%     20.0-against-21.4-per-cent agreement, is withdrawn entirely.
%     If you want the cost claim, put a std::chrono clock around the
%     weight computation alone on an idle machine. Ten lines, and 4.4.1
%     gets a measurement instead of an estimate.
%
%  --------------------------------------------------------------------
%  4. THE DAMPING CHECK NOW PASSES FOR ALL THREE, IDENTICALLY.
%         V1  predicted 6.20e8  observed 6.16e8   -0.62%
%         V2  predicted 4.53e8  observed 4.51e8   -0.62%
%         V3  predicted 4.53e8  observed 4.50e8   -0.62%
%     and solving for the exponent instead of assuming it returns 9.997,
%     9.997 and 9.997 against the 10 of equation (4.6). Three runs, the
%     same residual of 0.62 per cent in each, over nine orders of
%     magnitude. A systematic 0.62 per cent is almost certainly the
%     rounding of the two-figure 5.3e-18 I took from 4.7.1; with the exact
%     5.25935e-18 now in hand the agreement is as close as the data allow.
%
%  --------------------------------------------------------------------
%  5. STILL TO CHECK IN THE IMPLEMENTATION (unchanged from last time).
%     (a) The printed lambda and mu do not satisfy the printed equations
%         in absolute terms: with lambda = 2.33425 and Ne = 45.0219,
%         (4.5)-(4.6) give mu = 0.72 against the 3.24e-9 logged. The error
%         quantity mu uses is smaller than lambda's by a factor of 6.83,
%         the same at every operating point, so a constant is missing from
%         one of the two equations or they use different normalisations.
%     (b) Ne = 45.0219 occluded, 45.0095 nominal. Not a count of
%         measurements, though 4.2 and 4.7.1 write N_e as though it were.
%     (c) 4.7.1's "2.5 grey levels" should be 2.01: the lambda pair is
%         Variant 3's, 53 to 0.0407037, a ratio of 1302, so the intensity
%         error falls by 36.1 from 72.7.
% ======================================================================

\subsection{Convergence under occlusion}
\label{sec:occlusion}

The three variants are now compared on the occluded target. This is the
configuration for which the weighting was introduced, and the one in which the
three schemes are expected to separate. The reference image is acquired from
the unoccluded target, while the images acquired during servoing carry the
opaque elliptical region described in Section~\ref{sec:test-configurations},
so that every measurement falling inside it is in error and the error is
systematic rather than random: the occluder does not add noise to the cost, it
displaces the pose at which the cost is least. The poses, the filter bank and
the gain schedule are unchanged from Section~\ref{sec:nominal}, and the
occluder covers $7.4\%$ of the measurements at the initial pose and $5.80\%$
at the desired pose, the difference following from the retreat in depth
between them.

\begin{figure}[H]
    \centering

    \begin{minipage}[b]{0.30\textwidth}
        \centering
        \includegraphics[width=\textwidth]{Figs/C/fig1_a.PNG}

        \vspace{1mm}
        \textbf{(a)}
    \end{minipage}
    \hfill
    \begin{minipage}[b]{0.30\textwidth}
        \centering
        \includegraphics[width=\textwidth]{Figs/D/actual.PNG}

        \vspace{1mm}
        \textbf{(b)}
    \end{minipage}
    \hfill
    \begin{minipage}[b]{0.30\textwidth}
        \centering
        \includegraphics[width=\textwidth]{Figs/D/initiialerrur.PNG}

        \vspace{1mm}
        \textbf{(c)}
    \end{minipage}

    \vspace{5mm}

    \begin{minipage}[b]{0.30\textwidth}
        \centering
        \includegraphics[width=\textwidth]{Figs/D/V1.PNG}

        \vspace{1mm}
        \textbf{(d)}
    \end{minipage}
    \hspace{0.05\textwidth}
    \begin{minipage}[b]{0.30\textwidth}
        \centering
        \includegraphics[width=\textwidth]{Figs/D/V2V3.PNG}

        \vspace{1mm}
        \textbf{(e)}
    \end{minipage}

    \caption{Occlusion case, experimental setup: (a)~desired image
    $\mathbf{I}^{*}$, acquired from the unoccluded target; (b)~current image
    $\mathbf{I}$ at the initial pose, showing the occluder; (c)~initial image
    error $\mathbf{I}-\mathbf{I}^{*}$; (d)~final image error, Variant~1;
    (e)~final image error, Variant~3, that of Variant~2 being visually
    indistinguishable from it. In~(d) the residual extends over the whole
    image, the camera having settled at a pose displaced from the desired
    one; in~(e) it is confined to the occluder, which no scheme can remove
    because the reference contains no counterpart to it.}
    \label{fig:occ_setup}
\end{figure}

\begin{figure}[H]
    \centering

    \begin{minipage}[b]{0.30\textwidth}
        \centering
        \includegraphics[width=\textwidth]{Figs/D/camV12.PNG}

        \vspace{1mm}
        \textbf{(a)}
    \end{minipage}
    \hfill
    \begin{minipage}[b]{0.30\textwidth}
        \centering
        \includegraphics[width=\textwidth]{Figs/D/posV12.PNG}

        \vspace{1mm}
        \textbf{(b)}
    \end{minipage}
    \hfill
    \begin{minipage}[b]{0.30\textwidth}
        \centering
        \includegraphics[width=\textwidth]{Figs/D/errurV12.PNG}

        \vspace{1mm}
        \textbf{(c)}
    \end{minipage}

   \caption{Occlusion case, Variant~1 (no weighting, $\mathbf{D}=\mathbf{I}$):
   (a)~camera velocities; (b)~translational and rotational positioning errors;
   (c)~feature error $\lVert\mathbf{e}\rVert^{2}$.}

   \label{fig:occ_v1}

\end{figure}

\begin{figure}[H]
    \centering

    \begin{minipage}[b]{0.30\textwidth}
        \centering
        \includegraphics[width=\textwidth]{Figs/D/camV22.PNG}

        \vspace{1mm}
        \textbf{(a)}
    \end{minipage}
    \hfill
    \begin{minipage}[b]{0.30\textwidth}
        \centering
        \includegraphics[width=\textwidth]{Figs/D/posV22.PNG}

        \vspace{1mm}
        \textbf{(b)}
    \end{minipage}
    \hfill
    \begin{minipage}[b]{0.30\textwidth}
        \centering
        \includegraphics[width=\textwidth]{Figs/D/errurV22.PNG}

        \vspace{1mm}
        \textbf{(c)}
    \end{minipage}

   \caption{Occlusion case, Variant~2 (Tukey weighting): (a)~camera
   velocities; (b)~translational and rotational positioning errors;
   (c)~feature error $\lVert\mathbf{e}\rVert^{2}$.}

   \label{fig:occ_v2}

\end{figure}

\begin{figure}[H]
    \centering

    \begin{minipage}[b]{0.30\textwidth}
        \centering
        \includegraphics[width=\textwidth]{Figs/D/camV32.PNG}

        \vspace{1mm}
        \textbf{(a)}
    \end{minipage}
    \hfill
    \begin{minipage}[b]{0.30\textwidth}
        \centering
        \includegraphics[width=\textwidth]{Figs/D/posV32.PNG}

        \vspace{1mm}
        \textbf{(b)}
    \end{minipage}
    \hfill
    \begin{minipage}[b]{0.30\textwidth}
        \centering
        \includegraphics[width=\textwidth]{Figs/D/errurV32.PNG}

        \vspace{1mm}
        \textbf{(c)}
    \end{minipage}

   \caption{Occlusion case, Variant~3 (Hermite-informed weighting):
   (a)~camera velocities; (b)~translational and rotational positioning errors;
   (c)~feature error $\lVert\mathbf{e}\rVert^{2}$.}

   \label{fig:occ_v3}

\end{figure}
%%% All three captions said "feature-error norm ||e||". The plots and the
%%% text are of ||e||^2.

\paragraph{Results:}
Figures~\ref{fig:occ_v1}--\ref{fig:occ_v3} show the camera velocities, the
positioning error and the feature error against iteration, and
Table~\ref{tab:Occlusion_pose} collects the final pose errors. All three runs
converge in the sense of Section~\ref{sec:metrics}, the commanded velocity
having fallen below $10^{-5}\,\mathrm{m\,s^{-1}}$ in every case, and none
diverges or oscillates. The separation between the three is entirely a
separation in accuracy.

The separation is large. The unweighted scheme settles $1293\,\mu\mathrm{m}$
and $1.700^{\circ}$ from the desired pose, against $2.47\,\mu\mathrm{m}$ and
$0.0165^{\circ}$ for Tukey weighting and $3.16\,\mu\mathrm{m}$ and
$0.0152^{\circ}$ for the Hermite-informed weighting: a factor of between four
and five hundred in translation, and of about a hundred in rotation, between
the baseline and either weighted variant. Weighting an occluded scene repays,
many times over, the few tenths of a degree it costs on a clean one. The
contribution of robust weighting to a direct scheme is therefore not marginal
but decisive, and this is the principal quantitative result of the chapter.

\paragraph{The unweighted scheme minimises a displaced cost.}
The manner of the baseline's failure is more instructive than its magnitude,
and it is the clearest demonstration in this thesis of why a direct scheme
requires robustification. Variant~1 does not fail to converge, nor does it
come to rest at a higher value of the cost than the weighted variants. It
attains the \emph{smallest} final feature error of the three,
$1.5507\times10^{10}$ against $1.5734\times10^{10}$ for Variant~2 and
$1.5731\times10^{10}$ for Variant~3 --- a cost lower by $1.4\%$ --- at a pose
error larger by a factor of four hundred. The quantity compared is in each
case the unweighted residual, so the three are directly comparable, and the
ordering is unambiguous: the scheme that minimises the cost best is the scheme
that positions the camera worst.
%%% That the printed ||e||^2 is the unweighted residual follows from the
%%% numbers themselves: the occluder contributes a large share of the cost
%%% and Variants 2 and 3 reject it, so a weighted figure would be far
%%% SMALLER than Variant 1's. It is larger.

The occluded measurements have not raised the cost at the correct pose; they
have displaced the pose at which the cost is least, and Variant~1 locates that
displaced minimum accurately. Its pose, its feature error, its gain and its
damping are all unchanged to six significant figures between iterations~$400$
and~$600$: the run is not still converging towards the right answer, it has
finished converging on the wrong one. Two consequences follow. The first is
practical: no criterion computed from the same corrupted measurements that
drive the control can detect this failure, which is the argument of
Section~\ref{sec:metrics} met here as a measurement rather than asserted,
since a practitioner monitoring the residual would conclude that the
unweighted run was the most successful of the three. The second is that the
comparison vindicates the formulation of Section~\ref{sec:robust-framework} as
a problem of estimation rather than of optimisation: the difficulty is not
that the cost is hard to minimise, but that minimising it is the wrong thing
to do.

\paragraph{Neither weighted variant is uniformly the more accurate.}
Between Variants~2 and~3 the outcome is mixed. Variant~2 attains the smaller
translational error, $2.47\,\mu\mathrm{m}$ against $3.16\,\mu\mathrm{m}$, a
margin of $28\%$; Variant~3 attains the smaller rotational error,
$0.0152^{\circ}$ against $0.0165^{\circ}$, a margin of $8\%$. The
structure-informed weighting is therefore not shown by this experiment to
position more accurately than generic Tukey weighting under occlusion, and in
translation it is the less accurate of the two.
%%% This is the sentence you will be asked about in the viva. Having it in
%%% your own words, unhedged, is a far better position than having it put
%%% to you -- and it costs you nothing, because of the paragraph below.

The comparison should nevertheless be given its proper weight, which is
slight. Both weighted variants recover the desired pose to within three
micrometres and two hundredths of a degree, against $1293\,\mu\mathrm{m}$ and
$1.7^{\circ}$ for the baseline. A difference of seven tenths of a micrometre
between them is without practical consequence at that scale, rests on a single
run of each, and is not of a size that this experiment can establish as
systematic; the difference between either of them and the baseline is larger
by nearly three orders of magnitude and is not in doubt. What the occlusion
experiment establishes is the value of robust weighting in a direct scheme,
not a ranking among robust weightings.

One further observation is recorded, both because it is a genuine difference
between the two and because it qualifies the comparison above. At
iteration~$600$ the commanded velocity of Variant~3 is
$6.3\times10^{-6}\,\mathrm{m\,s^{-1}}$ against
$1.8\times10^{-7}\,\mathrm{m\,s^{-1}}$ for Variant~2, larger by a factor of
thirty-four, while its pose error is unchanged from iteration~$400$ to within
half a per cent. The residual velocity is therefore not descent that has yet
to be exhausted; it produces no net progress, and Variant~3 comes to rest with
a larger terminal jitter than Variant~2. This is consistent with the weights
being recomputed from the current image at every iteration, so that
$\mathbf{D}^{\mathrm{H}}$ continues to vary slightly along the trajectory even
after the pose has settled --- the time-varying behaviour discussed in
Section~\ref{sec:robust-stability}. The effect is small, it does not
destabilise the loop, and a single run cannot establish its cause; it is
reported because it is the one systematic difference the occlusion experiment
reveals between the two weighted variants, and because it plausibly accounts
for the translational margin above.
%%% Last time I offered two readings of Variant 3's larger velocity at
%%% iteration 400: unexhausted descent, or chatter. The 600 data decide it
%%% -- the pose did not improve, so it is the second. Reporting it is
%%% better than omitting it: it gives you a mechanism for the 28 per cent,
%%% and a mechanism is defensible where a bare number is not.

What the experiment does show in Variant~3's favour is that comparable
accuracy is reached by a different route, and the route matters for the reason
given in Section~\ref{sec:hermite-weight-limitation}: the structure-normalised
weighting suppresses measurements because they carry no usable image
structure, whereas the residual-based weighting suppresses them because they
disagree with the model, and only the first of these is a property of the
image rather than of the current pose estimate.
Section~\ref{sec:weight-maps} examines the mechanism directly. Establishing
any ordering on accuracy would require a study over several initial poses,
and a weight map for Variant~2, which
was not recorded.

\paragraph{A remark on the two configurations.}
The final errors of Table~\ref{tab:Occlusion_pose} are some two to three
orders of magnitude smaller than those of Table~\ref{tab:ch4-nominal}, and the
two tables must not be read as though occlusion had improved the positioning.
The explanation lies in the adaptive gain, and it is the mechanism identified
in Section~\ref{sec:nominal}. There the run was found to stall not because the
residual had been minimised but because the gain decayed with it:
$\lambda\propto\norm{\mathbf{e}_I}^{2}$ by~\eqref{eq:lambda}, and the
commanded velocity is proportional to the residual in addition, so
$\norm{\mathbf{V}_c}$ falls roughly as the cube of the error and the servo
stalls while a small misalignment remains. Under occlusion the intensity error
cannot decay to zero: the reference contains no counterpart to the occluder,
so a residual of fixed magnitude persists over some six per cent of the image
however well the camera is positioned. The gain is computed from that
intensity error and is held up by it, while the weighting removes the same
measurements from the descent direction. The servo therefore continues to
drive on the valid measurements at an undiminished gain, the cubic collapse of
the nominal case is replaced by a velocity proportional to the error alone,
and the run converges geometrically until it meets the limits of the
simulation rather than those of the gain schedule.

The measurements confirm this, and admit an independent check. At
iteration~$600$ the gain stands at $\lambda=2.334$, $2.193$ and $2.192$ in the
three occluded runs, against $\lambda=0.0407$ at the end of the nominal run: a
factor of between $54$ and $57$. The damping provides the check, being driven
by the same intensity error with a different exponent. Since
$\mu\propto\norm{\mathbf{e}_I}^{10}$ by~\eqref{eq:mu} while
$\lambda\propto\norm{\mathbf{e}_I}^{2}$, the ratio of the damping between the
two configurations must be the fifth power of the ratio of the gain. The
predicted ratios are $6.20\times10^{8}$, $4.53\times10^{8}$ and
$4.53\times10^{8}$; the observed are $6.16\times10^{8}$, $4.51\times10^{8}$
and $4.50\times10^{8}$, each low by $0.62\%$. Solving for the exponent instead
of assuming it returns $9.997$ in all three runs, against the value of $10$
in~\eqref{eq:mu}. Two quantities logged independently, separated by nine
orders of magnitude, agreeing across three runs to within two thirds of a per
cent: the account above is not an interpretation placed on the data but a
relation the data satisfy.

Three observations follow. First, the micrometre accuracies of
Table~\ref{tab:Occlusion_pose} are a property of the gain schedule under a
persistent residual and are not evidence that occlusion assists convergence;
the comparison that matters is between the three variants of that table, run
under identical conditions, and not between the two tables. Second, the same
mechanism explains why Variant~1 fails as sharply as it does rather than
merely drifting: it too enjoys the sustained gain, and it uses it to converge
efficiently upon the displaced minimum. Third, and more usefully, the stall
reported in Section~\ref{sec:nominal} is thereby shown to be a limitation of
the gain schedule inherited from the implementation and not of the weighting
or of the features, since holding $\lambda$ up --- by whatever means ---
improved the final accuracy by three orders of magnitude. A gain that does not
vanish with the error is accordingly the most immediate improvement available
to the scheme, and the present chapter provides, inadvertently, the evidence
for it.

\begin{table}[H]
    \centering
    \caption{Occlusion case: positioning error at iteration~600 for the three
    weighting variants, relative to the desired pose
    $\mathbf{t}^{*}=(0,0,1300)\,\mathrm{mm}$,
    $\boldsymbol{\omega}^{*}=(0,0,0)^{\circ}$, from the initial pose
    $\mathbf{t}_{0}=(40,-40,1100)\,\mathrm{mm}$,
    $\boldsymbol{\omega}_{0}=(12,-12,6)^{\circ}$. Translations in
    micrometres, rotations in degrees. The occluder covers $5.80\%$ of the
    measurements at the desired pose. One run per variant.}
  \label{tab:Occlusion_pose}
    \small
    \begin{tabular}{@{}lrrrrrrrr@{}}
        \toprule
        & \multicolumn{3}{c}{Translation ($\mu$m)}
        & \multicolumn{3}{c}{Rotation (deg)}
        & \multicolumn{2}{c}{Norm} \\
        \cmidrule(lr){2-4}\cmidrule(lr){5-7}\cmidrule(lr){8-9}
        Variant & $\Delta T_x$ & $\Delta T_y$ & $\Delta T_z$
                & $\Delta\omega_x$ & $\Delta\omega_y$ & $\Delta\omega_z$
                & $\norm{\Delta\mathbf{t}}$ ($\mu$m)
                & $\norm{\Delta\boldsymbol{\omega}}$ (deg) \\
        \midrule
        1 --- none    & $675.61$  & $30.796$ & $-1102.19$ & $-0.23572$  & $-1.68298$  & $-0.01961$  & $1293.1$ & $1.69952$ \\
        2 --- Tukey   & $-0.0181$ & $2.3040$ & $-0.8999$  & $+0.015123$ & $+0.006479$ & $-0.000543$ & $2.474$  & $0.016462$ \\
        3 --- Hermite & $1.8213$  & $1.7542$ & $-1.8881$  & $+0.014549$ & $+0.004326$ & $-0.000473$ & $3.156$  & $0.015187$ \\
        \bottomrule
    \end{tabular}
\end{table}

% ======================================================================
%  CONVERSION ARITHMETIC, from the three consoles at full precision.
%  Translations in metres x 1e6; rotations in radians x 180/pi.
%
%  V1  t = 0.000675613, 3.07961e-05, -0.00110219
%        -> 675.61, 30.796, -1102.19 um            |dt| = 1293.1 um
%      w = -0.00411405, -0.0293736, -0.000342328
%        -> -0.23572, -1.68298, -0.01961 deg       |dw| = 1.69952 deg
%  V2  t = -1.81389e-08, 2.304e-06, -8.99905e-07
%        -> -0.0181, 2.3040, -0.8999 um            |dt| = 2.474 um
%      w = 0.000263949, 0.000113084, -9.48219e-06
%        -> 0.015123, 0.006479, -0.000543 deg      |dw| = 0.016462 deg
%  V3  t = 1.82129e-06, 1.75417e-06, -1.88811e-06
%        -> 1.8213, 1.7542, -1.8881 um             |dt| = 3.156 um
%      w = 0.00025394, 7.5504e-05, -8.25116e-06
%        -> 0.014549, 0.004326, -0.000473 deg      |dw| = 0.015187 deg
%
%  WHAT MOVED BETWEEN 400 AND 600
%    V1  1293.1 -> 1293.1 um, 1.69952 -> 1.69952 deg. Frozen.
%    V2  2.909  -> 2.474  um (-15.0%) ; 0.016434 -> 0.016462 deg (+0.2%)
%    V3  3.139  -> 3.156  um (+0.5%)  ; 0.015203 -> 0.015187 deg (-0.1%)
%    |Tc| V1 1.43e-6 -> 2.16e-6 ; V2 2.40e-6 -> 1.83e-7 ; V3 2.83e-5 -> 6.26e-6
%
%  ONE ASYMMETRY WORTH NOTICING, not asserted in the text because a single
%  run will not support it. Variant 2's translational error is carried
%  almost entirely by Delta T_y and Delta T_z, with Delta T_x at 0.018 um,
%  three orders below them. Variant 3's three components are equal to
%  within four per cent, all about 1.8 um. An isotropic residual is what
%  one expects when nothing systematic remains; Variant 2's strongly
%  anisotropic one suggests a direction still weakly constrained, and it
%  is the direction in which it happens to have done well. Worth a line of
%  enquiry in the multi-pose study, not a claim here.
%
%  CONSEQUENT EDITS ELSEWHERE
%   (i) 4.5 says the trajectories of 4.7 "show no oscillation, no sign
%       reversal and no secondary minimum", offered as evidence that the
%       slowly-varying-weights assumption was not violated. Variant 3's
%       terminal jitter is a small counter-example -- not instability, but
%       not nothing either. Soften to "no oscillation of the trajectory,
%       no sign reversal and no secondary minimum was observed, the
%       terminal jitter reported in Section 4.7.2 aside", or add the
%       one-line check I suggested there: log ||D_k - D_{k-1}||_F against
%       the iteration. With the jitter now measured, that check would turn
%       4.5's assumption into a result.
%  (ii) 4.7.1: the lambda and mu quoted there are Variant 3's, exactly
%       0.0407037 and 5.25935e-18. Use the exact values -- the 0.62 per
%       cent residual in the damping check above is almost certainly the
%       rounding of the two-figure 5.3e-18.
% (iii) 4.7.1's "Gauss-Newton thereafter": the occluded mu is 2.4e-9, nine
%       orders above the nominal 5.3e-18. Check whether mu*diag(H) is
%       still negligible under occlusion before that sentence is allowed
%       to stand for both configurations.
%  (iv) Any earlier section quoting 1291 um, 2.91 um or 3.12 um for these
%       runs: the iteration-600 figures are 1293.1, 2.474 and 3.156.
%
%  LABELS USED ABOVE
%    sec:test-configurations  4.6.3      sec:metrics            4.6.4
%    sec:nominal              4.7.1      sec:weight-maps        4.7.3
%    sec:robust-framework     4.2        sec:robust-stability   4.5
%    sec:hermite-weight-limitation 4.4.4 sec:future             future work
%    tab:ch4-nominal          table of 4.7.1
%    eq:lambda , eq:mu        gain and damping
% ======================================================================
